\subsection{A criterion that targets the full simulation problem}
\label{tht:reveal}

Here is a precise criterion for the information that classical queries
must reveal. It avoids confusing coordinate derivatives with information
gained by observing a spin. For a law $\nu$, acceptance function $Q$, and
depth-$k$ classical query tree with transcript $\mathcal F$, write
\[
 v=\operatorname{Var}_{\nu}Q,
 \qquad G_{\nu}(Q,\mathcal F)
   =\operatorname{Var}_{\nu}\bigl(\mathbb E_{\nu}[Q\mid\mathcal F]\bigr).
\]
The law of total variance gives
$G_{\nu}(Q,\mathcal F)=v-\mathbb E_{\nu}\operatorname{Var}(Q\mid\mathcal F)$.
The tree may be adaptive; insisting on a fixed set of sites would be a
stronger, unnecessary condition.

\begin{proposition}[Polynomial reveal progress]
\label{tht:reveal-prop}
Let $\mathcal C$ be a class of input laws closed under conditioning any
set of coordinates to a positive-probability assignment. Suppose that
for every $\nu\in\mathcal C$, every $T$-query acceptance function $Q$,
and every $0<a\le1$, whenever $\operatorname{Var}_{\nu}Q\ge a$ there is
a transcript of depth at most $K(T,1/a)$ with
\[
 G_{\nu}(Q,\mathcal F)\ge g(T,a)>0.
\]
Assume $K(T,1/a)$ and $1/g(T,a)$ have polynomial upper bounds uniform
over the class. Then $\mathcal C$ has full AA, uniformly over its laws
and all such conditionings. Conversely, a uniform full-AA theorem for
$\mathcal C$ implies this progress property, with $g(T,a)=3a/4$ and a
polynomial $K$.
\end{proposition}

\begin{proof}
For sufficiency take $0<\epsilon<1$ and let $a=\epsilon^2/2$; larger
errors are trivial. At each transcript leaf whose
conditional variance is at least $a$, run a progress tree of depth at
most $K(T,1/a)$; freeze every other leaf. Restricted quantum algorithms
still use at most $T$ queries. Let $V_j$ be the expected conditional
variance after round $j$. It is nonincreasing and at most one. If
$V_j\ge2a$, the probability of active leaves is at least $a$: indeed
$V_j\le a+\Pr(\text{active})$, since each conditional variance is at
most one. The next round therefore decreases $V_j$ by at least
$a g(T,a)$. After
\[
 R=\left\lceil\frac{1}{a g(T,a)}\right\rceil+1
\]
rounds, $V_R<2a=\epsilon^2$; otherwise the total decrease would exceed
one. At each final leaf output the conditional expectation of $Q$.
Cauchy--Schwarz gives mean absolute error at most $\sqrt{V_R}<\epsilon$,
with depth $RK(T,1/a)$. Computing these expectations is permitted by
the query model.

Conversely, simulate $Q$ to mean absolute error $a/4$, fixing a random
seed by averaging if necessary, and let
$\mathcal F$ be the simulator's transcript. Clip its output $D$ to
$[0,1]$. Conditional expectation minimizes mean squared error, hence
\[
 \mathbb E\operatorname{Var}(Q\mid\mathcal F)
 \le\mathbb E(Q-D)^2\le\mathbb E|Q-D|\le a/4.
\]
If $v\ge a$, this gives $G\ge v-a/4\ge3a/4$. The simulator's query
bound is the required polynomial $K$.
\end{proof}

This is a characterization of \emph{hereditary} full AA, namely a
uniform theorem after arbitrary pinnings. It does not assert that
full AA for a single unconditioned family automatically extends to its
exponentially rare conditional laws. The characterization is an
operational target, not by itself a computable test on coupling matrices.

Even under a product law, one-site reveal variance is insufficient:
$Q(x)=(1+x_1x_2)/2$ has variance $1/4$ but
$\mathbb E[Q\mid X_i]=1/2$ for each individual coordinate. Two queries
reveal $Q$ exactly. Conversely, Curie--Weiss magnetization has tiny
flip derivatives in the ordered phase but can be estimated by sampling
coordinates. These two examples explain why transferring a
flip-influence inequality is neither the necessary nor the most direct
route to the primary target.

For classical Gibbs laws, pinning preserves interaction order and the
sum of absolute incident interaction strengths: terms meeting fixed
spins become lower-order terms. A candidate high-temperature theorem
should therefore control reveal progress uniformly over the induced
fields and boundary conditions. For measured noncommuting quantum
Gibbs states, conditioning measurement outcomes need not equal taking
the Gibbs state of the naively pinned Hamiltonian; that additional
step must be justified separately.
